\section*{Author notation, parameter regime and Section 3}
$I,J,A,B,S$ denote vertex-sets and $T$ edge-sets. $E(S):=\binom S2$. $G$ denotes a simple graph on vertex-set $[n]$. ``$T\subseteq E([n])$'' is omitted in summation when there is no confusion. $y(n)=O(x(n))$ means there is some absolute constant $c$ such that $y(n)\le cx(n)$ for all $n$.
Throughout the paper,
\[\epsilon=\text{any positive parameter (wlog }\epsilon<1/40),\quad\omega=n^{1/2-4\epsilon},\quad\tau=\frac\epsilon{200}\frac{\log n}{\log\log n},\quad d=\frac\epsilon{100}\tau.\]
\subsection*{3.1 Non-exact case: a new perspective}
Given a graph $G$ we can think of a degree-$d$ pseudo-expectation as assigning a number $\widetilde{\mathbb E}x_S$ to each subset $S\subseteq[n]$ of size $\le d$, so that the resulting vector $\widetilde{\mathbb E}x$ looks indistinguishable to the expectation resulted from the case when a random-$\omega$ clique is planted, from the view of degree-$d$ SOS. To make up such an assignment we first go beyond to slightly larger subsets of size $\tau$. We define an ``approximate distribution'' on size $\le\tau$-cliques in $G$ then use it to generate pseudo-expectation on all size $\le d$-subsets.
\subsubsection*{3.1.1 $\zeta$-function and M\"obius inversion}
Given $n$-vertex graph $G$, let $p(G)\in\mathbb R^{2^{[n]}}$ be the max-clique-indicator vector, then $q(G):=\zeta p(G)$ is a vector supported exactly on all cliques in $G$, where $\zeta$ is the $2^{[n]}\times2^{[n]}$ matrix
\[\zeta(A,B)=1\text{ iff }A\subseteq B,\qquad\forall A,B\subseteq[n].\tag{3.1}\]
In particular, if $G$ itself is a single clique then $q(G)$ is the clique-indicator. $\zeta_{a,b}$ denotes the submatrix of $\zeta$ on rows $\binom{[n]}{\le a}$ and columns $\binom{[n]}{\le b}$; use similar notation on all related vectors. Consider the plant-situation where $G$ is indeed a single random clique. Suppose its distribution is represented by a plant-distribution vector $p_{\mathrm{plant}}\in\mathbb R^{2^{[n]}}$. Let the output-expectation $q_{\mathrm{out}}$ be indicator-vector of cliques in $G$ in expectation. Then
\[q_{\mathrm{out}}=\zeta p_{\mathrm{plant}}.\tag{3.2}\]
We call such a pair $(p_{\mathrm{plant}},q_{\mathrm{out}})$ a plant-setting.
\paragraph{Definition 3.1 (Two plant-settings).} The exact plant-setting $(p_0,q_0)$ is:
\[p_0(S)=\binom n\omega^{-1}\text{ if }|S|=\omega\text{ and }0\text{ otherwise},\tag{3.3}\]
and
\[q_0(S)=(\zeta p_0)(S)=\frac{\binom{n-|S|}{\omega-|S|}}{\binom n\omega}.\tag{3.4}\]
I.e. a random size-$\omega$ subset is chosen to be the planted clique. The independent plant-setting $(p_1,q_1)$ is:
\[p_1(S)=(\omega/n)^{|S|}(1-\omega/n)^{n-|S|}\quad\forall S\subseteq[n],\tag{3.5}\]
and
\[q_1(S)=(\zeta p_1)(S)=(\omega/n)^{|S|}.\tag{3.6}\]
I.e. any vertex is included in the planted clique w.p. $\omega/n$ independently.
Thus $\zeta$ reveals the basic linear relations between $(p_{\mathrm{plant}},q_{\mathrm{out}})$. It is upper-triangular (row- and column-indices ordered in a size-ascending way), invertible, and the inverse is the M\"obius inversion matrix:
\[\zeta^{-1}(A,B)=(-1)^{|B\setminus A|}\text{ if }A\subseteq B,\text{ and }0\text{ otherwise}.\]
Note $(\zeta_{a,a})^{-1}=(\zeta^{-1})_{a,a}$ for all $a\le n$. Moreover, if let the pseudo-expectation be defined as $\widetilde{\mathbb E}x=p\in\mathbb R^{2^{[n]}}$ for some vector $p$, then the ``full'' $2^{[n]}\times2^{[n]}$ moment matrix is
\[M_{\mathrm{SOS}}=\zeta\operatorname{diag}(p)\zeta^\top.\tag{3.7}\]
In particular, if $p$ is a nonnegative vector then $M_{\mathrm{SOS}}$ is immediately PSD.
\subsubsection*{3.1.2 The non-exact pseudo-expectation}
Given any $G$, first construct a degree-$\tau$ ``approximate plant-distribution'' $p_\tau(G)$, which simulates the plant-distribution (Definition3.1) in the sense that they give similar output-expectations. Require $p_\tau(G)$ to be supported on size $\le\tau$-cliques in $G$. Then take $\widetilde{\mathbb E}x=\zeta_{d,\tau}p_\tau(G)$ so the result inherits the linear structure posed by $\zeta$. From approximation it seems taking $\zeta_{\tau,\tau}^{-1}(q_1)_\tau$ would suffice, while to make it supported on cliques, same as in [13] add a clique-indicator factor:
\[p_\tau(G)(S)=2^{|\binom S2|}\mathrm{Cl}_S(G)\left(\zeta_{\tau,\tau}^{-1}(q_1)_\tau\right)(S)\quad\forall S\subseteq[n],\ |S|\le\tau.\tag{3.8}\]
Here $\mathrm{Cl}_S(\cdot)$ is the clique indicator function and $2^{|\binom S2|}$ is for re-normalization.
\paragraph{Definition 3.2.} $\forall S\subseteq[n]$, the normalized clique-indicator is function \[\widetilde{\mathrm{Cl}}_S(G):=2^{|\binom S2|}\mathrm{Cl}_S(G).\tag{3.9}\]
$\widetilde{\mathrm{Cl}}(G)$ denotes the column vector over a family of $S$'s, always clear from context.
\paragraph{Definition 3.3.} The non-exact pseudo-expectation is
\[\widetilde{\mathbb E}_{\mathrm{nonexact}}=\zeta_{d,\tau}p_\tau(G)=\zeta_{d,\tau}\left(\widetilde{\mathrm{Cl}}(G)\circ\zeta_{\tau,\tau}^{-1}\right)(q_1)_\tau\in\mathbb R^{\binom{[n]}{\le d}},\tag{3.10}\]
where $\circ$ is Hadamard product.\footnote{In general $(M_1\circ M_2)M_3\ne M_1\circ(M_2M_3)$, but they are equal if $M_1$ is a column vector.} In short, this refines the construction in [13] by one step: factor through size-$\tau$ subsets (in the only non-trivial way) so the size-$d$ output inherits linear relations posed by $\zeta$. The resulting moment matrix is
\[M_{\mathrm{nonexact}}(G)=\zeta_{d/2,\tau}\operatorname{diag}(p_\tau(G))(\zeta_{d/2,\tau})^\top,\tag{3.11}\]
similarly as (3.7).
\paragraph{Remark 3.4.} If $p_\tau(G)$ were nonnegative then PSDness would be immediate. Alas, this is not true by computation. PSDness could still possibly hold because $\zeta_{d/2,\tau}$ in (3.11) is degenerate.
\paragraph{Lemma 3.5 (Theorem1.5(1)).} For all $S\subseteq[n]$ with $|S|\le d$,
\[\widetilde{\mathbb E}_{\mathrm{nonexact}}x_S=\sum_{T:|V(T)\cup S|\le\tau}(\omega/n)^{|V(T)\cup S|}\chi_T.\tag{3.12}\]
\begin{proof}Note $\widetilde{\mathrm{Cl}}_S=\sum_{T\subseteq E(S)}\chi_T$ for all $S$. Now for $S,S'$ with appropriate size bound,
\[\left(\widetilde{\mathrm{Cl}}\circ\zeta_{\tau,\tau}^{-1}\right)(S,S')=\begin{cases}\sum_{T\in E(S)}\chi_T(-1)^{|S'\setminus S|},&S\subseteq S',\\0,&\text{o.w.}\end{cases}\]
Then
\begin{align*}&\left(\zeta_{d,\tau}(\widetilde{\mathrm{Cl}}\circ\zeta_{\tau,\tau}^{-1})\right)(S,S')\\&=\sum_{S'':S\subseteq S''\subseteq S'}\left(\sum_{T\subseteq E(S'')}\chi_T(-1)^{|S'\setminus S''|}\right)\\&=\sum_{T:V(T)\cup S\subseteq S'}\chi_T\left(\sum_{S'':V(T)\cup S\subseteq S''\subseteq S'}(-1)^{|S'\setminus S''|}\right)\\&=\sum_{T:V(T)\cup S\subseteq S'}\chi_T\delta_{S'=V(T)\cup S}\\&=\sum_{T:V(T)\cup S=S'}\chi_T.\end{align*}
Therefore,
\begin{align*}\widetilde{\mathbb E}_{\mathrm{nonexact}}x_S&=\left(\zeta_{d,\tau}(\widetilde{\mathrm{Cl}}\circ\zeta_{\tau,\tau}^{-1})(q_1)_\tau\right)(S)\\&=\sum_{S':|S'|\le\tau}\left(\sum_{T:V(T)\cup S=S'}\chi_T(\omega/n)^{|S'|}\right)\\&=\sum_{T:|V(T)\cup S|\le\tau}\chi_T(\omega/n)^{|V(T)\cup S|},\end{align*}
for all $S$ with $|S|\le d$.\end{proof}
\subsection*{3.2 The exact case}
In this subsection, we give a generic way to generate possible pseudo-expectations that satisfy Size Constraints (1.4). Define $\widetilde{\mathbb E}x_S$ in a top-down fashion: fix it for all $|S|=d$ first, then recursively set
\[\widetilde{\mathbb E}x_S\leftarrow\frac1{\omega-|S|}\sum_{i\notin S}\widetilde{\mathbb E}x_{S\cup\{i\}}\tag{3.13}\]
for smaller-sized $S$'s. Denote by $\widetilde{\mathbb E}_dx$ the assignments for $|S|=d$.
\paragraph{Definition 3.6.} The $d$-filtration matrix $\mathrm{Fil}_{d,=d}$, of dimension $\binom{[n]}{\le d}\times\binom{[n]}d$, is:
\[\mathrm{Fil}_{d,=d}(A,B)=\begin{cases}\binom{\omega-|A|}{d-|A|}^{-1},&A\subseteq B\ (|B|=d),\\0,&\text{otherwise}.\end{cases}\tag{3.14}\]
\paragraph{Definition 3.7.} Given vector $\widetilde{\mathbb E}_dx$ assigning a value to each $d$-subset $S\subseteq[n]$, the exact pseudo-expectation it generates is
\[\widetilde{\mathbb E}x:=\mathrm{Fil}_{d,=d}\widetilde{\mathbb E}_dx.\tag{3.15}\]
\paragraph{Lemma 3.8.} The pseudo-expectation in Definition3.7 satisfies the Size Constraints (1.4), regardless of the choice of $\widetilde{\mathbb E}_dx$.
\begin{proof}For any $S\in\binom{[n]}{<d}$, take $v_S\in\mathbb R^{\binom{[n]}{\le d}}$,
\[v_S(S')=\begin{cases}\omega-|S|,&S'=S,\\-1,&S'\supseteq S,\ |S'\setminus S|=1,\\0,&\text{otherwise}.
\end{cases}\]
It suffices to show $v_S^\top\mathrm{Fil}_{d,=d}=0$. But this is a direct check.\end{proof}
The generated $\widetilde{\mathbb E}$ should further satisfy Clique, PSDness and Default Constraints. Default is not a problem as long as $\widetilde{\mathbb E}x_\emptyset>0$, since we can rescale everything by $(\widetilde{\mathbb E}x_\emptyset)^{-1}$ without affecting other constraints.
\subsection*{3.3 The exact pseudo-expectation}
Take the pseudo-expectation for $|S|=d$ in the form
\[\widetilde{\mathbb E}x_S=\sum_{T:|V(T)\cup S|\le\tau}\chi_TF(|V(T)\cup S|),\]
for a function $F$ called a $d$-generating function. For general $|S|\le d$, (3.14) gives, putting $u:=d-|S|$,
\begin{align*}\widetilde{\mathbb E}x_S&=\frac1{\binom{\omega-d+u}u}\sum_{T:|V(T)\cup S|\le\tau}\chi_T\left[\sum_{c=0}^u\binom{|V(T)\cup S|-d+u}c\right.\\&\qquad\left.\cdot\binom{n-|V(T)\cup S|}{u-c}F(|V(T)\cup S|+u-c)\right].\tag{3.16}\end{align*}
\paragraph{Lemma 3.10.} Any exact pseudo-expectation generated by (3.16) satisfies Clique and Size Constraints (1.3),(1.4).
\begin{proof}For Clique Constraints, note (3.16) only depends on $|V(T)\cup S|$, so by grouping terms $\widetilde{\mathbb E}x_S=\sum_{T:|V(T)\cup I\cup J|\le\tau}M(I,J;T)\chi_T$ factors through $\widetilde{\mathrm{Cl}}_{I\cup J}=\sum_{T\subseteq E(I\cup J)}\chi_T$. I.e., $M(I,J)(G)=0$ if $\widetilde{\mathrm{Cl}}_{I\cup J}(G)=0$. It satisfies Size Constraints by Lemma3.8.\end{proof}
\paragraph{Definition 3.11 (Exact $d$-generating function).}
\[F(x):=\frac{(x+8\tau^2)!}{(8\tau^2)!}(\omega/n)^x.\]
\paragraph{Remark 3.12.} The design of $F$, especially its first factor, is technical and the ultimate goal is to make the resulting $M$ positive. The numerator $(x+8\tau^2)!$ will be used in Proposition6.28, where $8\tau^2$ can be replaced by larger polynomials in $\tau$. The $(8\tau^2)!$ denominator is just for convenience (Remark3.14).
\paragraph{Definition 3.13.} The exact moment matrix $\widetilde M$ is
\[\widetilde M(A,B)=\sum_{T:|V(T)\cup A\cup B|\le\tau}\widetilde M(A,B;T)\chi_T\quad\forall |A|,|B|\le d/2,\]
where, putting $u=d-|A\cup B|$,
\begin{align*}\widetilde M(A,B;T)&=\frac1{\binom{\omega-d+u}u}\left[\sum_{c=0}^u\binom{|V(T)\cup A\cup B|-(d-u)}c\right.\\&\qquad\cdot\binom{n-|V(T)\cup A\cup B|}{u-c}\\&\qquad\left.\cdot\frac{(|V(T)\cup A\cup B|+u-c+8\tau^2)!}{(8\tau^2)!}(\omega/n)^{|(V(T)\cup A\cup B)|+u-c}\right].\tag{3.17}\end{align*}
